\section{Characterization of the neural-state reader}
\label{sec:reader}
The reader is an auxiliary interface component. It compresses downstream firing rates into 16 continuous prefix tokens for a frozen Qwen3-0.6B model. A two-layer projector with a 1,024-unit hidden layer is trained to generate structured Russian captions containing stimuli, an odor label, dopamine input, and thresholded behavioral readouts. The direct-input comparison uses the same projector design with a vector of stimulus identities, sides, rate categories, odor label, and dopamine condition.

We corrected a data-splitting error before the evaluation reported here. Canonicalization and deduplication reduce 6,234 stored rows to 3,906 unique state--caption examples. All examples sharing a downstream state are assigned to one partition, including examples with conflicting captions. Any state group containing sugar and bitter is reserved for the interaction evaluation. The remaining split has 2,723 training, 578 validation, and 578 test examples; 27 examples form the interaction set. Feature frequency, normalization, and odor vocabulary are fitted using training data only. The audit finds no overlap between partitions in downstream-state identity, state--caption identity, or the final selected feature vectors.

We trained three initializations on this fixed split and selected each checkpoint by validation loss. Test negative log-likelihood (NLL) is the mean negative log-probability per target token; lower values are better. It is evaluated on all 578 test examples after training. Generation metrics use the same seeded random subset of 100 test examples for every initialization. Table~\ref{tab:reader} reports the mean and sample standard deviation across the three initializations, not uncertainty over independent datasets.
\begin{table}[t]
\caption{Reader characterization on the corrected split. NLL uses 578 test rows; stimulus-set accuracy uses the same 100 generated captions; interaction accuracy uses 27 sugar--bitter rows. Values are mean $\pm$ sample standard deviation over three initializations of one fixed split. Interaction majority-class accuracy is 81.5\%. Shuffled-state generation was not evaluated.}
\label{tab:reader}
\centering\small
\begin{tabular}{@{}lrrr@{}}
\toprule
Input & Test NLL $\downarrow$ & Stimuli exact (\%) & Interaction (\%) \\
\midrule
Neural-state reader & 0.135 $\pm$ 0.022 & 56.7 $\pm$ 3.1 & 55.6 $\pm$ 23.1 \\
Direct stimulus input & 0.043 $\pm$ 0.004 & 89.0 $\pm$ 2.6 & 16.0 $\pm$ 2.1 \\
Reader, shuffled states & 0.737 $\pm$ 0.031 & -- & -- \\
\bottomrule
\end{tabular}
\end{table}


The trained reader uses state information: shuffling states at test time raises caption loss in all three runs. However, the direct-input model gives lower caption loss and higher exact stimulus-set accuracy. The held-out sugar--bitter probe is a separate limitation: reader proboscis accuracy ranges from 37.0\% to 81.5\%, while predicting the most frequent class reaches 81.5\%. The tested reader therefore does not provide a reliable account of this interaction. These observations characterize this trained interface; they do not imply a general limit on neural-to-language decoding.

This dataset predates the current chemical encoder. Its word-odor inputs were produced by a text-embedding-based projection into selected olfactory populations. The present scores must not be interpreted as reader performance on measured DoOR chemical responses, new memory states, or the case-study dialogue distribution. The narrator can still describe those experiments from explicit records, with that information access disclosed.

A separate physiological language experiment uses synthetic event histories and Qwen3-4B (Appendix~\ref{app:physiology-pilot}). Its follow-up holds the physiological states fixed and compares two decoder interfaces over three paired initializations (Appendix~\ref{app:question-decoder}). These experiments address a different task and representation; their results are reported separately from the caption reader above.
